\subsubsection{BIT point update + prefix sum}

\paragraph{Idea intuitiva.}
El \textbf{Fenwick Tree} (BIT) almacena sums parciales de manera que \texttt{sum(i)} (suma de $[1, i]$) y \texttt{add(i, delta)} cuestan $O(\log n)$. La representacion es ingeniosa: el nodo $i$ cubre el rango $[i - \text{lowbit}(i) + 1, i]$. Asi, \texttt{sum(i)} itera \texttt{i -= lowbit(i)}, sumando los nodos que cubren un prefijo creciente, y \texttt{add(i, delta)} itera \texttt{i += lowbit(i)}, propagando el cambio a todos los nodos cuyo rango contiene a $i$. La idea: cada indice $i$ se descompone en una suma de bloques disjuntos de tamano potencias de 2 (su ``lowbit representation'').

\paragraph{Propiedades clave (invariantes).}
\begin{itemize}\setlength\itemsep{0pt}
	\item Indice \textbf{1-based} (no usar $i=0$).
	\item \texttt{lowbit(i) = i \& (-i)}.
	\item Cada nodo cubre exactamente un rango de tamano \texttt{lowbit(i)}.
	\item La operacion debe ser \textbf{invertible} (suma, xor, no min/max).
\end{itemize}

\paragraph{Codigo clave.}
Las dos operaciones basicas:
\begin{verbatim}
void add(int i, long long d) {
    for (; i <= n; i += i & -i) bit[i] += d;
}
long long sum(int i) {
    long long s = 0;
    for (; i > 0; i -= i & -i) s += bit[i];
    return s;
}
\end{verbatim}

\paragraph{Complejidad.}
\begin{itemize}\setlength\itemsep{0pt}
	\item $O(\log n)$ por operacion. Memoria $O(n)$.
	\item Constante muy pequena ($\sim 1$): en la practica, $\sim 3x$ mas rapido que segment tree.
\end{itemize}

\paragraph{Donde tocar para modificar.}
\begin{itemize}\setlength\itemsep{0pt}
	\item \textbf{Suma -> max/min/xor}: cambiar el \texttt{+=} por la operacion (con identidad apropriada). \textbf{No} funciona para \texttt{min} de forma natural (no es invertible).
	\item \textbf{Tamano dinamico}: aniadir un metodo \texttt{resize(newSize)} que preserve valores.
	\item \textbf{Tamano grande}: con $n = 10^7$, mejor comprimir coordenadas.
	\item \textbf{Operador no conmutativo}: usar segment tree en su lugar.
\end{itemize}

\paragraph{Trampas y errores frecuentes.}
\begin{itemize}\setlength\itemsep{0pt}
	\item Indices 0 vs 1: el BIT es \textbf{1-based}. Si el problema da indices 0-based, aniadir \texttt{+1} en cada llamada.
	\item El metodo \texttt{sum} itera \texttt{i -= i \& -i}; aniadir o quitar 1 rompe el invariante.
	\item Si el tipo es \texttt{int}, posibles overflows; usar \texttt{long long}.
	\item Para rango $[l, r]$, el query es \texttt{sum(r) - sum(l-1)} (cuidado con $l = 1$).
\end{itemize}

\paragraph{Explicacion linea por linea.}
\textbf{Estructura del BIT:}
\begin{itemize}\setlength\itemsep{0pt}
	\item \verb|struct BIT { int n; vector<long long> t; };| -- encapsula el Fenwick Tree.
	\item \verb|BIT(int n): n(n), t(n+1, 0) {}| -- constructor: tamano $n+1$ para indexacion 1-based.
	\item \verb|void add(int i, long long delta)| -- agrega \texttt{delta} en la posicion \texttt{i}.
	\item \verb|for(; i <= n; i += i & -i) t[i] += delta;| -- avanza por los ancestros de \texttt{i} (\texttt{lowbit}), sumando en cada uno.
	\item \verb|long long sum(int i)| -- calcula la suma de \texttt{a[1]} a \texttt{a[i]}.
	\item \verb|for(; i > 0; i -= i & -i) r += t[i];| -- desciende removiendo el bit menos significativo.
	\item \verb|long long rangeSum(int l, int r)| -- retorna la suma en $[l, r]$ mediante \texttt{sum(r) - sum(l - 1)}.
\end{itemize}
\textbf{Programa principal:}
\begin{itemize}\setlength\itemsep{0pt}
	\item \verb|BIT bit(10);| -- crea BIT de tamano 10.
	\item \verb|bit.add(3, 5);| -- suma 5 en la posicion 3.
	\item \verb|bit.add(7, 2);| -- suma 2 en la posicion 7.
	\item \verb|cout << bit.sum(5) << "\n";| -- imprime suma de $[1, 5]$ (que es 5).
	\item \verb|cout << bit.rangeSum(1, 10) << "\n";| -- imprime suma total (7).
\end{itemize}

\paragraph{Codigo.}
Ver \texttt{cpp.tex}, seccion \textit{BIT point update + prefix sum}.

\vspace{0.5em}

